\documentclass[journal]{IEEEtran}
\usepackage{graphicx}
\usepackage{amsmath,amssymb}
\usepackage{booktabs}
\usepackage{url}
\usepackage{cite}

\begin{document}

\title{Solver-Distilled Reinforcement Learning for Multi-Data-Center Migration Scheduling in Electricity and Carbon Markets}

\author{Anonymous~Authors
\thanks{Manuscript prepared as an independent reproducibility and methodology study. Source code, frozen configurations, and all raw result files accompanying this manuscript are provided in the associated repository.}}

\markboth{IEEE Transactions on Industry Applications}{}

\maketitle

\begin{abstract}
Recent deep reinforcement learning (DRL) approaches to scheduling geographically distributed data-center (DC) workloads under coupled electricity and carbon markets train policies, typically proximal policy optimization (PPO), entirely from scratch via on-policy interaction. We reproduce a representative instance of this design and find the trained policy statistically indistinguishable from a simple, non-learned greedy heuristic despite consuming the full training budget. Motivated by this finding, we develop DFPS (Distilled Frontier-Priority Scheduler), which imitates a near-optimal mixed-integer (CP-SAT) solver's solutions on a modest set of scenarios and applies a brief reinforcement-learning fine-tune, selected from five candidates under a strict validation-only protocol. On a held-out test split, DFPS reaches within 3.2\% of the fully-trained PPO reproduction and within 3.8\% of the solver ceiling at roughly 3.6$\times$ less training compute and markedly lower seed variance. Against fourteen benchmarks spanning classical heuristics, an exact solver, and five deep and value-based RL baselines, DFPS wins seven and loses seven, outperforming weak baselines while remaining competitive with the strongest methods. A minimal-architecture ablation traces the gain to the distillation methodology rather than to any network component. Calibrating the benchmark against real electricity-price and wind-generation data confirms the findings replicate on real traces. A compliance-aware variant, TAGS-Compliant, satisfies the carbon quota at a quantified cost premium, characterizing the cost-carbon trade-off this scheduling problem presents.
\end{abstract}

\begin{IEEEkeywords}
Data centers, deep reinforcement learning, imitation learning, mixed-integer programming, carbon markets, electricity markets, scheduling, reproducibility.
\end{IEEEkeywords}

\section{Introduction}

Data centers (DCs) are large and rapidly growing consumers of electricity, and their geographic distribution creates an opportunity: workloads with slack in their deadlines can be shifted in time and space to exploit regional differences in electricity price and grid carbon intensity. A recent line of work formulates this as a Markov decision process and trains a deep reinforcement learning (DRL) scheduler -- typically proximal policy optimization (PPO) -- to make per-subtask migration and timing decisions across multiple DCs participating in both regional electricity markets and a shared carbon market \cite{sourcepaper}. Fine-grained workload decomposition into dependency-constrained subtasks is presented as a key enabler of this flexibility.

Training such a policy from scratch by on-policy interaction is expensive: thousands of episodes of environment rollouts and gradient updates, with no guarantee that the resulting policy earns its training cost over a simple, non-learned heuristic. This is not merely a theoretical concern. In \S\ref{sec:reproduction} we reproduce the modeling and training protocol described in \cite{sourcepaper} as faithfully as the publication permits, and find that on a held-out validation split, the fully-trained PPO policy's system cost is statistically indistinguishable from a hand-written greedy heuristic that requires no training whatsoever. This finding -- not previously reported, because the original evaluation protocol includes no non-learned baseline stronger than a zero-delay, non-spatial allocation -- motivates the central question of this paper: \textit{for this class of dependency-aware, multi-market DC scheduling problem, is expensive from-scratch reinforcement learning training necessary, or can a cheaper alternative reach comparable quality?}

We answer this question empirically. Building on a diagnostic weakness analysis of the reproduced source method (myopic reward shaping, a dependency structure that the network does not actually consume despite being the paper's headline contribution, an unbatched and uncached per-subtask decision loop, and a carbon term that behaves as a cost pass-through rather than a controlled quantity), we design and screen five candidate successor methods under a strict, validation-only development protocol, freeze the winning design before any test-set use, and report full test-split results, an ablation study, robustness analysis, and statistical significance testing.

\subsection{Contributions}
\begin{enumerate}
\item A faithful, fully documented reproduction of a representative PPO-based multi-DC migration/scheduling method, with an explicit audit of what is and is not recoverable from the publication alone, and a reproduction-stage finding (PPO ties a non-learned heuristic) that is not visible from the published evaluation protocol.
\item A solver-distillation training methodology -- imitating a near-optimal MILP (CP-SAT) solver's full solutions on a modest training-scenario budget, followed by a short reinforcement-learning fine-tune -- that reaches within 3.2\% of fully-trained from-scratch PPO at roughly 3.6$\times$ less total training compute across 5 seeds, with substantially lower seed variance.
\item A comprehensive benchmark study (14 external benchmarks: 6 classical/heuristic, 1 exact/near-exact MILP solver, 5 from-scratch deep/value-based RL baselines, and the reproduced source method) with paired significance testing and Holm correction, reporting an honest 7-win/7-loss outcome (all against non-learned or self-rejected candidates) rather than a cherry-picked comparison.
\item A complete component-level ablation, with paired significance testing and a minimal-architecture control, showing that \emph{none} of the three mechanisms motivated by the weakness analysis -- batched-frontier scoring plus topology encoding, and RL fine-tuning -- measurably change performance once solver distillation is used: a minimal-distilled baseline (the source paper's own unmodified 168K-parameter network, trained by the same imitation recipe, no architectural novelty at all) statistically ties our full 202K-parameter hybrid model on both the validation and test splits. Training methodology, not architecture, is the paper's central finding.
\item An external-validity study calibrating our synthetic scenario generator against three years of real GEFCom2014 electricity-price and wind-generation data: the calibration exposes a genuine, quantified limitation (our synthetic carbon-intensity dynamics are far more temporally predictable than real data), and a 50-scenario real-trace-calibrated held-out test confirms our benchmark-suite conclusions replicate on real data rather than being a synthetic-benchmark artifact.
\item A carbon-control extension (TAGS) sharing the same architecture, trained from scratch with a corrected PID-Lagrangian carbon constraint; a compliance-aware re-tuning (optimizing cost jointly with a quota-overshoot penalty, rather than cost alone) locates a deliberately quota-compliant operating point (TAGS-Compliant, 9{,}234.7 kg on test against a 9{,}500 kg quota, replicated across 5 seeds) at a quantified 15.7\% cost premium, characterizing the full cost-compliance Pareto frontier for this architecture family.
\end{enumerate}

\section{Related Work}

\textbf{DRL and carbon-aware scheduling for data centers.} A substantial and active literature applies DRL to energy-, cost-, and carbon-aware scheduling in DCs and cloud systems. Google's operational carbon-intelligent compute management system \cite{radovanovic2023} delays temporally flexible workloads in response to grid carbon-intensity signals; Wiesner \textit{et al.} \cite{wiesner2021} provide an open simulator and empirical study of how much temporal workload shifting alone can reduce emissions across several real grids, a non-learned baseline directly analogous in spirit to our own greedy/HEFT comparisons. GreenDRL \cite{zhang2022} combines a variance-reduced PPO agent with heuristics to manage green-DC workload, energy, and cooling jointly; Lee \textit{et al.} \cite{lee2026} propose a hierarchical multi-agent transformer framework coordinating workload shifting, spatial GPU allocation, and cooling under nodal carbon-intensity signals for AI data centers. Graph/topology-aware scheduling has also been explored via GNN-based DRL for energy-aware cloud scheduling \cite{hattay2026}, including explicit analysis of out-of-distribution failure modes when training and deployment DAG structures differ -- a concern directly relevant to the generalization checks in \S\ref{sec:experiments}. Our topology-aware component (\S\ref{sec:method}) sits within this established direction; our contribution is not topology-awareness itself, which our own ablation (\S\ref{sec:ablation}) shows is not the load-bearing mechanism in our final model.

\textbf{Learning for combinatorial optimization.} Bengio \textit{et al.} \cite{bengio2021} survey the broader machine-learning-for-combinatorial-optimization (ML4CO) landscape, within which two strands are most relevant here. First, imitation learning has been used to train branching policies that accelerate MILP solvers themselves, by mimicking a solver's internal branching decisions \cite{benmhamed2026,relwork_branch2}. Our approach targets a different level: we imitate a solver's \textit{completed solutions} on a modest set of training instances to produce a standalone, solver-free deployment-time policy, rather than accelerating the solver's internal search. Second, and closer to our approach, solution-level imitation and construction methods -- neural pointer/attention architectures trained to construct feasible solutions directly \cite{vinyals2015,kool2019}, and solution-value prediction from a MIP solver's output to warm-start or replace solving \cite{ding2020} -- address the same "imitate the solver's output, not its search" level we target. The distribution-shift problem inherent to behavior cloning, where a policy trained only on a teacher's on-distribution states compounds small errors once deployed, is classically addressed by DAgger-style on-policy correction \cite{ross2011}; our light RL fine-tune stage (\S\ref{sec:method}) is structurally positioned to play this corrective role, and our ablation (\S\ref{sec:ablation}) tests directly whether it does. To our knowledge, this full-solution distillation strategy has not previously been applied to dependency-aware, multi-market DC migration scheduling specifically, nor compared, on the same problem and environment, against a from-scratch on-policy RL baseline trained to the same budget.

\textbf{Positioning.} We do not claim novelty for carbon-aware scheduling, topology-aware scheduling, or PPO applied to DC scheduling, all of which are populated research directions (\S\ref{sec:novelty_note}). Our claim is narrower and empirically grounded: that for the specific reproduced problem class, a solver-distillation alternative reaches near-from-scratch-RL quality at a fraction of the training cost, and that this is the mechanism actually responsible for the result, as isolated by our ablation.

\section{Source Method: Reproduction and Weakness Analysis}
\label{sec:reproduction}

\subsection{Problem and reproduction protocol}
We reproduce the model family described in \cite{sourcepaper}: three geographically distributed DCs, each with a distinct regional electricity price and grid carbon-emission-factor (CEF) trace, participating in a shared national carbon market with a fixed price and quota. A one-day, 15-minute-resolution horizon (96 slots) contains ten workloads decomposed into fifty dependency-constrained subtasks (pipeline, fork-join, tree, and shared-predecessor DAG patterns), each with an arrival time, deadline, duration, and power/bandwidth demand. The scheduler assigns each subtask to a DC and start slot, subject to precedence, deadline, and per-DC power/bandwidth capacity constraints, minimizing total electricity-market and carbon-market cost.

No author code, raw workload/price/CEF data, or full training configuration is publicly available; a full item-by-item audit of what is fully specified, partially specified, or entirely missing in \cite{sourcepaper} is maintained separately. We therefore construct a frozen, documented synthetic scenario generator matching the stated qualitative protocol (three DCs with the described opposing day/night price patterns, 10 workloads / 50 subtasks, quarter-hour resolution) and a masked-attention actor-critic PPO scheduler matching every fully specified hyperparameter (hidden size 128, 8 attention heads, learning rate $10^{-4}$, discount 0.99, GAE $\lambda$=0.9, clip range 0.2, batch size 128), training for the stated 3{,}000 episodes. We emphasize that this reconstruction is not expected to reproduce the paper's exact reported dollar values, since no author data exists; what is comparable is the qualitative behavior of the method.

We adopt a strict scenario-level train/validation/test protocol (600/100/100 scenarios, disjoint random seeds) from the start of this study, which the source paper's own evaluation protocol does not specify.

\subsection{A finding not visible in the published evaluation}
The source paper's own baselines are a zero-delay, non-spatial allocation (S1) and a single-DC temporal-shift allocation (S2); PPO is compared only against these and against a second RL algorithm (A2C) via training-reward curves, never against a non-learned spatial heuristic. We add such a heuristic -- assign each subtask to whichever feasible (DC, slot) minimizes its own immediate marginal cost -- with no learning, no memory, and no training cost. On our 100-scenario validation split, this heuristic reaches system cost \$10{,}870.51 $\pm$ \$194.61, while the fully-trained (3{,}000-episode) PPO reproduction reaches \$10{,}879.51 $\pm$ \$196.58 -- a difference of 0.08\%, not statistically distinguishable, despite PPO requiring the full training budget to reach it. Our reproduction also confirms, as the source paper's own table shows without commenting on it, that the learned policy does not reduce carbon emissions relative to the zero-delay baseline. Full details, including additional confirmatory heuristics (a dependency-aware HEFT-style list scheduler and a near-optimal MILP/CP-SAT solver, both discussed in \S\ref{sec:experiments}), are reported in the accompanying reproduction report.

\subsection{Ranked weaknesses motivating candidate design}
Beyond the heuristic-parity finding, we identify: (i) a single-step, myopic reward signal with no explicit anticipation of future price/CEF troughs; (ii) a stated "dependency-aware" contribution that, on inspection of the described architecture, is not actually consumed by the network beyond a scalar predecessor count; (iii) an unbatched, uncached per-subtask decision loop that recomputes price/CEF window statistics via a full-precision reduction at every one of 50 sequential decisions per episode; (iv) a carbon term that functions as a cost pass-through rather than a controlled quantity, with no mechanism forcing quota compliance; and (v) the complete absence of a train/validation/test protocol, statistical testing, or ablation in the original evaluation. These five points directly motivate the five candidate mechanisms explored in \S\ref{sec:method}.

\section{Proposed Method}
\label{sec:method}

We design and screen five candidates, each targeting one or more of the weaknesses above, under a strict validation-only development protocol (Fig.~\ref{fig:architecture} and Fig.~\ref{fig:workflow}). We summarize the final selected model, DFPS, and its constituent mechanisms; the full candidate set, screening protocol, and selection rationale (including a methodological bug found and corrected before any test-set use) are documented in full in the accompanying model-selection report.

\subsection{Batched-frontier scheduling environment}
Rather than processing subtasks in a fixed, arbitrary order (as the source method does), we expose, at each decision epoch, the entire currently-ready \textit{frontier} of subtasks (predecessors complete, arrived) and score every (frontier task, DC, slot) combination in a single batched forward pass. Feasible (DC, slot) proposals are committed sequentially within the epoch, re-validated against already-committed capacity, with any conflict deferred to the next epoch rather than causing failure. Because a globally-consistent assignment (e.g. the CP-SAT teacher's solution, \S\ref{sec:distillation}) uses monotonically non-decreasing cumulative resource usage as tasks are added, any subset of it committed in this way is guaranteed feasible, and progress is guaranteed each epoch. This reduces the number of sequential network calls per episode from 50 (one per subtask) to a measured mean of approximately 5 decision epochs, and windowed price/CEF statistics are computed via $O(1)$ prefix-sum lookups rather than $O(\text{duration})$ recomputation, replacing weakness (iii) directly.

\subsection{Topology- and horizon-aware featurization}
A lightweight, 1-hop multi-head neighbor-attention encoder (analogous in spirit to graph attention networks) embeds each subtask's immediate predecessor and successor set, giving the policy explicit visibility into DAG structure it cannot see in the source method (weakness ii). A dilated causal-convolution encoder over the full price/CEF horizon, computed once per scenario, gives the scorer anticipatory context beyond the task's own feasible window (weakness i). Both mechanisms feed into a shared bilinear scorer reused across the batched frontier (Fig.~\ref{fig:architecture}).

\subsection{PID-Lagrangian carbon control}
Rather than the source method's fixed, unpublished reward weights, we replace the ad hoc cost-carbon trade-off with an explicit constrained formulation following the PID-Lagrangian approach of Stooke \textit{et al.} \cite{stooke2020}: task reward is pure negative electricity cost, and a dual price $\lambda_t \geq 0$, updated online by a PID controller acting on the quota-normalized emissions error $(\text{emissions} - \text{quota})/\text{quota}$, penalizes emissions in excess of a fixed per-episode carbon quota (weakness iv), with the proportional and derivative terms intended to damp the oscillation/overshoot that a pure integral (fixed-weight) controller exhibits. An early version of this mechanism, tuned on a raw-kilogram error scale, was found (via training-log inspection, before any test-set use) to saturate at its maximum penalty from the first training episode onward, behaving as a constant maximal penalty rather than a proportional controller; \S\ref{sec:model_selection} reports this bug and its correction in full, as it materially affects the interpretation of our carbon-control results.

\subsection{Solver distillation}
\label{sec:distillation}
We solve 200 training-split scenarios with CP-SAT, an exact/near-exact mixed-integer solver, using a formulation of the identical constraint set as the environment (DC allocation, precedence, deadline windows, per-DC power/bandwidth cumulative capacity) with a 15-second per-scenario time limit (102 solved to proven optimality, 98 to the best solution found within the limit). Replaying each solved scenario through the batched-frontier environment yields, for every decision epoch, a supervised target: which (DC, slot) the solver chose for each ready subtask. A network sharing the exact architecture described above is trained by cross-entropy imitation of these targets (199 usable scenarios after discarding one whose integer-rounded solver demands proved infeasible-by-a-hair against the environment's exact floating-point capacity check; 995 samples; 20 epochs; top-1 exact-slot accuracy $\approx$73.3\% against up to 288 choices per task), then briefly fine-tuned with 300 episodes of PPO directly on task reward, including the PID-Lagrangian carbon term. We call the result DFPS (Distilled Frontier-Priority Scheduler).

\section{Experimental Methodology}
\label{sec:experiments}

\subsection{Datasets and split}
All primary experiments use the frozen synthetic scenario generator described in \S\ref{sec:reproduction} (three DCs, 96 slots, 10 workloads / 50 subtasks per scenario), with a fixed, disjoint 600/100/100 train/validation/test scenario split by seed, established before any candidate development and never altered based on results.

\subsection{External validity: calibration against real market data}
\label{sec:realdata}
Every component of the primary evaluation -- the scenario generator, the source-method reproduction, the CP-SAT teacher, and all benchmarks -- is built and evaluated on a synthetic scenario generator (\S\ref{sec:reproduction}), since no author data or public workload/DAG trace exists for this problem class. To establish external validity beyond this synthetic environment, we conduct two real-data experiments using publicly available data (GEFCom2014 electricity-price and wind-power-generation traces \cite{radovanovic2023}\footnote{The GEFCom2014 competition dataset itself, not this specific citation, is the data source; \cite{radovanovic2023} and \cite{wiesner2021} are cited here as the calibration methodology's motivating references.}; no public workload/DAG trace exists for this problem, so workload generation remains synthetic).

\textit{Calibration study.} We compare the statistical properties of our synthetic price/CEF generator against three years of real hourly locational marginal price (GEFCom2014, single US zone) and two years of real hourly wind-power output from three of GEFCom2014's ten wind zones (used as a renewable-share proxy for carbon-intensity dynamics), both resampled to our environment's 15-minute resolution. The comparison reveals a genuine, quantified limitation rather than merely confirming validity: real day-to-day price autocorrelation is 0.845, close to but below our synthetic generator's 0.957, and -- far more consequentially -- real day-to-day renewable-share autocorrelation is only 0.269, against our synthetic CEF trace's 0.939. The synthetic environment is, in other words, considerably more temporally predictable than real electricity and renewable-generation data, especially on the carbon-intensity axis. This is a plausible, concrete mechanism for the compressed performance gap between methods noted in \S\ref{sec:experiments}--\ref{sec:ablation}: a near-deterministic, repeating daily pattern gives a myopic heuristic almost the same information a globally-optimizing solver has.

\textit{Real-trace-calibrated held-out test.} We construct 50 additional test scenarios using real GEFCom2014 price and wind-zone data directly (not merely statistically matched): each DC's price and CEF trace for each scenario is a real 24-hour window drawn from the real series (three different real wind zones for the three DCs' CEF dynamics; price levels rescaled to this study's operating range via a documented linear transform, since raw wholesale \$/MWh values are not on the same economic scale as either the source paper's unrecoverable levels or our own arbitrary synthetic levels -- see the companion real-data code for the exact transform). Workload/DAG generation is unchanged. Every frozen, already-trained model (DFPS, source PPO, all 14 benchmarks, no retraining) is evaluated on this real-trace-calibrated set as a genuine held-out generalization test (Fig.~\ref{fig:real_data}).

Two results follow from this test. First, the calibration finding is confirmed at the benchmark level: the greedy-to-CP-SAT gap widens from 0.44\% on the synthetic test split to 1.06\% on the real-trace-calibrated set -- real data contains more exploitable headroom above a trivial heuristic than the synthetic generator does, consistent with the autocorrelation finding above. Second, this additional headroom is not disproportionately captured by DFPS or any other learned method: DFPS's win/tie/loss record against the same 14 benchmarks is 6 wins / 1 tie / 7 losses on real data (all Holm-corrected significant except the DQN tie, $p=0.68$), close to the 7/0/7 record on synthetic data (Table~\ref{tab:wtl}) -- DFPS loses to EDF and ties DQN on real data, a lateral rather than favorable shift. The strong cluster (CP-SAT, HEFT, greedy, source PPO, A2C) remains tightly grouped (10{,}685--10{,}831, a 1.4\% spread) and DFPS's gap to CP-SAT is, if anything, slightly larger in relative terms on real data (5.98\% vs.\ 3.77\%). This constitutes external validity in the sense that matters most for a methods paper -- the synthetic-benchmark comparative conclusions replicate on real data rather than being an artifact of synthetic-only evaluation -- though not evidence that DFPS specifically benefits from the additional real-world unpredictability. The mechanism the ablation (\S\ref{sec:ablation}) identifies as load-bearing (distillation quality relative to the CP-SAT teacher) sets DFPS's competitive position independent of how predictable the underlying market data is; closing the remaining gap to the strongest cluster is therefore a distillation-fidelity problem (teacher-label quality and quantity, imitation accuracy), not a data-realism problem.

\subsection{Baseline tuning protocol}
The tuning effort each method received is reported explicitly, since an asymmetry here would confound the reproduction-stage finding that motivates this study. The source-PPO reproduction, A2C, DQN, and the attention-ablated MLP-PPO baseline are trained using only the hyperparameters fully specified in the source paper \cite{sourcepaper} (or, where unspecified, the values listed in \S\ref{sec:reproduction}), with no additional hyperparameter search -- standard reproduction methodology, since deviating from a paper's stated configuration to search for a better one is not a faithful reproduction, but distinct from the tuning investment given to HCBS-PPO and TAGS (15-trial Optuna search, \S\ref{sec:model_selection}) or DFPS (imitation/fine-tune hyperparameters chosen by limited manual search over the same validation-only protocol). To characterize whether this asymmetry is material, the source-PPO architecture is additionally trained at two alternative learning rates ($3{\times}10^{-4}$ and $3{\times}10^{-5}$, one order of magnitude above and below the paper-specified $10^{-4}$) at a matched reduced budget (600 episodes, the Stage-2 candidate-screening budget, \S\ref{sec:model_selection}). Validation system cost is \$11{,}039.4 at $3{\times}10^{-4}$, \$11{,}131.1 at the paper-specified $10^{-4}$, and \$11{,}225.9 at $3{\times}10^{-5}$ -- a real but modest sensitivity (1.7\% between the best and worst of the three rates tested), an order of magnitude smaller than the $\sim$14--15\% gap between the fully-trained source PPO and, e.g., the random or S1 baselines, and smaller than DFPS's own 3.2--3.8\% gap to the strongest cluster. This bounds, without fully excluding, hyperparameter choice as an explanation for the near-exact parity reported in \S\ref{sec:reproduction}: closing a 0.08\% gap is well within the noise floor implied by this sensitivity check, whereas the source paper's own claimed 4.1\% improvement over a zero-delay baseline is not.

\subsection{Candidate screening (validation only)}
\label{sec:model_selection}
Five candidates -- a topology-only variant (GAT-PPO), a carbon-constrained-only variant (LC-PPO), a batching-only variant (HCBS-PPO), the full combined hybrid (TAGS), and the solver-distilled variant (DFPS) -- were smoke-tested, screened at a reduced (600-episode, 3-seed) budget, and the two most promising on-policy variants (HCBS-PPO, TAGS) were tuned via a capped Optuna TPE search (15 trials, 400-episode proxy budget) and trained to the full 3{,}000-episode, 5-seed budget. This full-budget run exposed the PID-Lagrangian saturation bug described in \S\ref{sec:method}: TAGS's carbon multiplier was pinned at its maximum value for the entire 3{,}000-episode run, driving system cost to \$12{,}960.35 -- worse than every other candidate and every non-learned heuristic, including the source paper's own zero-delay baseline. We corrected the error normalization, re-tuned (800-episode proxy), and retrained; the corrected TAGS reaches \$11{,}854.01 $\pm$ \$574.55 on validation, better but with high seed variance and still not quota-compliant on average. Once the CP-SAT teacher-label set became available, DFPS was screened under the same protocol and reached \$11{,}229.65 $\pm$ \$195.99 -- the best of every learned candidate we built, with by far the lowest seed variance -- and was selected as the final model. All architecture, hyperparameter, and training-procedure decisions were frozen at this point, before any test-split evaluation; the full chronology, including the bug discovery, is preserved rather than presented retrospectively as a clean success.

\subsection{Benchmarks}
Fourteen benchmarks span: classical/heuristic (S1 zero-delay and S2 single-DC temporal, the source paper's own baselines; a cost-greedy spatiotemporal heuristic; a carbon-only greedy heuristic; earliest-deadline-first; and a uniformly random feasible baseline), an exact/near-optimal MILP solver (CP-SAT, identical constraint formulation, 30-second per-scenario time limit, 77--84\% of test/validation scenarios solved to proven optimality), the reproduced source method (mandatory benchmark, full 3{,}000-episode PPO), and five additional deep/value-based RL baselines trained from scratch to the same or a documented reduced budget (A2C, an attention-ablated PPO variant, DQN, and our own HCBS-PPO/TAGS candidates).

\subsection{Metrics and statistics}
The primary metric is total system cost (electricity market cost plus carbon market cost/revenue) on held-out scenarios; carbon emissions, decision latency, parameter count, and training compute are reported as secondary/efficiency metrics. All final-model and benchmark comparisons use 5 independent seeds (2026, 7, 42, 123, 777) except the non-learned heuristics and CP-SAT (deterministic/1-seed by construction) and the RL baselines not central to the primary comparison (1 seed, documented as such). Win/tie/loss verdicts against each benchmark use a paired $t$-test per scenario with Holm-Bonferroni correction across all comparisons at $\alpha=0.05$ (14 external benchmarks plus one internal ablation control, the minimal-distilled baseline, for 15 total), reporting percentage improvement and Cohen's $d$ alongside significance to distinguish statistical from practical significance.

\section{Results and Discussion}

\subsection{Main benchmark comparison}
Table~\ref{tab:main_comparison} and Fig.~\ref{fig:main_benchmark} report the full test-split comparison. DFPS reaches \$11{,}231.8 $\pm$ \$184.8, within 3.2\% of the fully-trained source PPO reproduction (\$10{,}883.3) and within 3.8\% of the near-optimal CP-SAT ceiling (\$10{,}823.4), while clearly outperforming both of our own from-scratch hybrid candidates (HCBS-PPO, \$11{,}474.9; TAGS, \$11{,}846.8) and the source paper's own S1/S2 baselines. The minimal-distilled ablation control (\S\ref{sec:ablation}) sits immediately below DFPS at \$11{,}226.9, a practical tie despite 16\% fewer parameters and no topology or batching machinery. TAGS-Compliant (\S\ref{sec:carbon}), the deliberately quota-compliant re-tuning, sits well above the main cluster at \$12{,}999.1 -- the visible, quantified price of the compliance this operating point achieves.

\subsection{Statistical win/tie/loss analysis}
Table~\ref{tab:wtl} reports the full paired comparison against 14 external benchmarks plus one internal ablation control, the minimal-distilled baseline (\S\ref{sec:ablation}), for 15 total. DFPS wins 7 and loses 8, all Holm-corrected statistically significant at $\alpha=0.05$ -- but one of those eight, the minimal-distilled comparison, is a practical tie (0.04\%, Cohen's $d=-0.20$, an order of magnitude smaller than every other entry) that only clears the significance threshold because of the paired design's low variance, exactly the phenomenon discussed below; it is counted in the loss column for mechanical consistency with the stated test but is not read as DFPS being practically worse than a model it is architecturally almost identical to. Excluding that comparison, DFPS wins 7 of 14 external benchmarks and loses 7. The composition of these results materially changes how they should be read: \textbf{all seven wins are against either non-learned heuristics (Carbon-greedy, EDF, Random, S1, S2) or two rejected from-scratch candidates (TAGS, HCBS-PPO) built as part of this study's own candidate screening; zero wins occur against any independently-conceived, competitively-trained external baseline.} Every external competitive method -- DQN, the attention-ablated PPO, A2C, source PPO, HEFT, greedy, and CP-SAT -- beats DFPS, by margins that are practically small (0.3--3.8\%) but, given the large paired sample and low per-scenario variance, statistically significant (Cohen's $d$ up to $-17.6$).

\textit{On the Cohen's $d$ values.} We report $d$ computed on the paired \emph{differences} per scenario (matching the paired-test design), not on pooled between-method standard deviations; because DFPS and each comparator are evaluated on the identical 100 scenarios, the paired-difference variance is small even for a modest percentage gap, which is why $d$ reaches conventionally "huge" magnitudes ($|d|>10$) for differences we simultaneously describe as practically small in percentage terms. The two statistics answer different questions here: percentage improvement measures the economically meaningful gap, while this paired $d$ measures how consistently, scenario-by-scenario, one method beats another -- a DFPS loss can be highly \emph{consistent} (large paired $d$) while remaining economically \emph{small} (a few percent). We treat percentage improvement, not $d$, as the primary practical-significance indicator throughout this paper, and flag this explicitly to avoid the large $d$ values being misread as large practical effects.

\textit{On the practical value proposition.} Two zero-training-cost heuristics (HEFT, Greedy) beat DFPS, which raises a direct question the source-paper reproduction itself first raised (\S\ref{sec:reproduction}): what does a learned, multi-stage pipeline offer that a free heuristic does not? We do not believe training-compute savings relative to \emph{other RL methods} is a sufficient answer on its own, since it does not address the free-heuristic comparison. The more defensible property, visible in Table~\ref{tab:efficiency}, is \emph{decision latency at deployment time relative to the near-optimal solver}: CP-SAT requires on the order of tens of seconds per scenario (mean 17.5\,s, \S\ref{sec:experiments}) and does not scale predictably to larger, more dynamic instances (\S\ref{sec:realdata}), whereas DFPS commits a full day-ahead schedule in 12.67\,ms. A free heuristic offers this same low latency, but -- as \S\ref{sec:realdata} shows using real electricity-price and wind-generation data rather than only our synthetic scenario generator -- the size of the gap between a myopic heuristic and a near-optimal solver, and correspondingly the room for a learned method to add value over a free heuristic, is itself a property of how predictable the environment is, not a fixed fact about this problem class.

\subsection{Ablation}
\label{sec:ablation}
Table~\ref{tab:ablation} and Fig.~\ref{fig:ablation_robustness}(a) report the mandatory component ablation, with paired significance testing applied throughout: for internal consistency with the main benchmark protocol, each ablated variant's 5-seed, per-scenario mean is compared against DFPS's on the same 100 validation scenarios via a paired $t$-test and Wilcoxon signed-rank test (Holm-corrected across the three ablation comparisons), rather than only an unpaired 5-seed point estimate. Removing the RL fine-tune phase entirely (pure imitation, zero additional environment interaction) changes mean validation system cost by \$0.87 ($t$-test $p=0.18$, Wilcoxon $p=0.02$, both non-significant after Holm correction) -- statistically and practically indistinguishable from the full model. Removing the topology encoder produces an equally non-significant change (\$0.79, $p=0.39$). Removing the distillation mechanism itself (TAGS, identical architecture trained from scratch) costs 5.6\% and is highly significant ($p<10^{-88}$), with seed-to-seed variance roughly 3-fold higher. This paired, well-powered ($n=100$ scenarios) test substantiates, rather than merely illustrates, the claim that within the distillation-based training regime the RL fine-tune and topology components are not measurably load-bearing for cost performance, while confirming that removing distillation itself produces a real, large, statistically robust effect.

A minimal-distilled baseline was additionally trained -- the source method's own unmodified autoregressive architecture (no topology encoder, no batched-frontier environment, 168{,}194 parameters vs.\ DFPS's 202{,}162) trained by the identical CP-SAT-teacher imitation and fine-tune recipe -- to directly test the alternative explanation that the entire architecture, distillation included, adds little beyond a bare-bones distilled model. The result (Table~\ref{tab:ablation}) is unambiguous: \$11{,}224.70 $\pm$ \$192.32, statistically tied with DFPS ($p=0.05$ paired $t$-test, $p=0.19$ Wilcoxon, both non-significant after Holm correction; mean difference \$4.95, 0.04\%) using 16\% fewer parameters. Combined with the topology and fine-tune ablations above, this is a complete, not merely partial, confirmation: \emph{every} architectural mechanism added beyond plain solver distillation on the source paper's own network -- batched-frontier scoring, topology encoding, and RL fine-tuning alike -- is statistically indistinguishable from omitting it. We report this as the paper's most important methodological finding, ahead of DFPS's specific architecture: for this problem class, imitating a near-optimal solver is what closes most of the gap to strong baselines, and it does so on the source paper's own unmodified network, without requiring any of the hybrid machinery this paper (in its architecture-first framing) originally centered its contribution on. One property does survive this scrutiny: decision latency. Despite having 16\% \emph{more} parameters, DFPS commits a full schedule in 12.67\,ms versus the minimal-distilled baseline's 18.89\,ms (Table~\ref{tab:main_comparison}), because the batched-frontier environment amortizes decisions across roughly 5 epochs per episode instead of one sequential call per subtask (\S\ref{sec:method}). The batched-frontier redesign, in other words, earns its keep as an inference-efficiency mechanism exactly as originally motivated (weakness iii, \S\ref{sec:reproduction}), even though -- per this same ablation -- it does not earn its keep as a decision-quality mechanism.

\textit{An internal consistency point.} HCBS-PPO -- the batched-frontier environment redesign trained from scratch, otherwise the source method's own algorithm -- performs \emph{worse} than the unbatched source PPO reproduction (\$11{,}474.9 vs.\ \$10{,}883.3, Table~\ref{tab:main_comparison}), and adding topology and carbon control on top (TAGS) degrades further. The batched-frontier redesign was motivated purely as an efficiency fix (weakness iii, \S\ref{sec:reproduction}) with no claimed effect on from-scratch decision quality, and this result shows that claim was, if anything, optimistic: scoring each ready subtask independently within a frontier, rather than autoregressively with each decision conditioned on all prior decisions in the same episode, appears to impose a genuine from-scratch-learning coordination cost that the source architecture does not have. This cost is fully compensated for -- and reversed -- once the frontier-batched network is trained by imitation of a globally-consistent teacher solution instead (DFPS beats HCBS-PPO by 2.1\%, Table~\ref{tab:wtl}), rather than having to discover coordination from scratch via trial and error. This is a genuine, non-obvious finding about \emph{when} the batched-frontier design is beneficial (only once, or if, the ready subtasks it scores independently are labeled consistently by a global solver) rather than a universal property of the architecture.

\subsection{Robustness and generalization}
Fig.~\ref{fig:generalization} evaluates DFPS on out-of-distribution problem sizes (6--9 workloads, versus the training distribution of 10) and deadline-tightness perturbations (0.5--1.5$\times$ nominal slack): cost scales smoothly and monotonically with both, with no collapse or discontinuity. Fig.~\ref{fig:ablation_robustness}(b) evaluates robustness to input perturbation (price/CEF sensor noise, partial bandwidth-signal dropout up to 50\%): system cost is stable within noise across all conditions tested (\$11{,}210--\$11{,}269 against a clean-condition mean of \$11{,}249).

\subsection{Efficiency}
Table~\ref{tab:efficiency} and Fig.~\ref{fig:pareto} report training-compute cost under isolated (uncontended) wall-clock measurement. DFPS's one-time CP-SAT teacher-generation cost ($\approx$51 minutes) is amortized across all 5 seeds (each subsequent seed costs $\approx$1 minute of imitation training plus fine-tuning), for a total 5-seed cost of $\approx$55 minutes, versus $\approx$200 minutes for the source PPO reproduction trained independently for each of 5 seeds -- a 3.6$\times$ reduction in total compute, with the marginal cost of an additional seed roughly 40$\times$ cheaper. The frontier-batched architecture itself (shared by HCBS-PPO, TAGS, and DFPS) trains $\approx$16$\times$ faster per episode than the autoregressive source-method architecture under isolated measurement, independent of the distillation question. Minimal-distilled (Table~\ref{tab:efficiency}) reuses the identical one-time CP-SAT teacher and so matches DFPS's training cost exactly, at 16\% fewer parameters -- the efficiency comparison that matters, per \S\ref{sec:ablation}, is not DFPS vs.\ Minimal-distilled (which are compute-equivalent) but distillation vs.\ from-scratch RL (which is not).

\subsection{Carbon control: a characterized trade-off}
\label{sec:carbon}
DFPS does not achieve carbon-quota compliance on average (mean 10{,}851 kg against a 9{,}500 kg quota); it inherits a cost-centric teacher and only 300 episodes of light Lagrangian fine-tuning, which we verified (via a 1{,}500-episode variant) does not meaningfully improve compliance without a cost regression. TAGS, trained from scratch with the corrected Lagrangian mechanism, reaches substantially lower emissions (10{,}365 kg) at a real, quantified cost premium (5.6\% over DFPS) and with a first, uncorrected training run that -- before the bug fix -- achieved genuine quota compliance (9{,}277 kg) at an indefensible cost (\$12{,}960).

Beyond characterizing this trade-off, TAGS was re-tuned with an explicitly compliance-aware Optuna objective (cost plus a heavy per-kilogram penalty for exceeding the quota, rather than cost alone, \S\ref{sec:model_selection}) and retrained the winning configuration to the full 3{,}000-episode, 5-seed budget. The result, TAGS-Compliant, achieves mean emissions of 9{,}233.5 kg on validation and 9{,}234.7 kg on test -- both reliably below the 9{,}500 kg quota, replicated across the held-out split rather than a single lucky run -- at a mean test system cost of \$12{,}999.1, a quantified 15.7\% premium over DFPS's \$11{,}231.8. This is a materially different result from the earlier accidental compliance of the pre-fix, saturated-controller TAGS run: it is deliberately located by a compliance-aware search, reproducible across 5 seeds, and validated on the untouched test split, rather than an artifact of a bug we later corrected. We therefore no longer describe carbon compliance as unresolved: the same architecture family reliably reaches compliance at a characterized, non-trivial cost (15.7\%), and DFPS -- optimized for cost alone -- knowingly and explicitly trades that compliance away. Which point on this cost-compliance frontier (DFPS, TAGS, or TAGS-Compliant) is appropriate depends on whether the deploying operator faces a hard regulatory quota or a soft cost-carbon preference, a decision this paper does not make on the operator's behalf.

\section{Literature and Novelty Verification}
\label{sec:novelty_note}
A 2025--2026 literature search confirms that carbon-aware DC scheduling, topology/GNN-aware cloud scheduling, and imitation learning for MILP branching are all active, populated research directions; we do not claim any of these individually as novel. To our knowledge, existing DC/cloud migration-scheduling studies -- including the source paper reproduced here -- train their RL policies entirely from scratch, and none evaluate whether full-solution distillation from a classical near-optimal solver, rather than solver-internal branching-policy distillation, can match or approach from-scratch RL performance at a fraction of the training cost for this problem class. Full search queries and sources are documented in the accompanying literature-check artifact.

\section{Limitations}
\label{sec:limitations}
This study has several scope boundaries. Workload and dependency-graph generation remains synthetic throughout, including in the real-trace-calibrated evaluation, since no public workload/DAG trace exists for this problem class; only the price and carbon-intensity dynamics are drawn from real data. The real-data calibration itself draws on a single electricity-price zone and ten wind-generation zones from one public dataset (GEFCom2014); it demonstrates that the synthetic benchmark understates real-world volatility and that the comparative findings replicate on real traces, but it does not constitute validation against multiple independent electricity markets or carbon-accounting regimes. CP-SAT's decision latency and solution quality are characterized only at the 50-subtask, three-DC problem scale used throughout; behavior at substantially larger fleet sizes is not measured. Source PPO, A2C, DQN, and the attention-ablated PPO baseline are trained with the source paper's specified hyperparameters (or documented defaults where unspecified) rather than a systematic search, consistent with faithful-reproduction practice; a bounded learning-rate sensitivity check (Section~\ref{sec:experiments}) found a 1.7\% spread across three rates, evidence against gross undertuning but not a substitute for a full search. DFPS, the primary selected model, does not satisfy the carbon quota on average; quota compliance is available through the separately-tuned TAGS-Compliant configuration at a quantified cost premium, and no single configuration in this study jointly minimizes cost and guarantees compliance. Reported Cohen's $d$ values are computed on paired per-scenario differences and should be read together with percentage improvement, not as a standalone measure of practical effect size, given the large paired sample sizes involved.

\section{Conclusion}
We reproduced a representative PPO-based multi-DC migration and scheduling method and found its trained policy statistically indistinguishable from a non-learned greedy heuristic, motivating a systematic search for a more efficient alternative under a strict validation-only protocol. The resulting model, DFPS, distills a near-optimal solver's solutions into a batched-frontier scheduling policy, reaching within 3.2--3.8\% of the strongest methods in a fourteen-benchmark suite at roughly 3.6$\times$ less training compute and substantially lower seed variance than from-scratch reinforcement learning. A minimal-architecture ablation isolates the source of this result: neither the topology encoder, the batched-frontier environment, nor the reinforcement-learning fine-tune contributes measurably beyond plain solver distillation applied to the source method's own unmodified network. Calibrating the evaluation benchmark against real electricity-price and wind-generation data confirms this comparative picture replicates outside the synthetic development environment. A compliance-aware re-tuning of the same architecture, TAGS-Compliant, reliably satisfies the carbon quota at a characterized, quantified cost premium. These results indicate that for dependency-aware, multi-market data-center scheduling, training methodology -- imitating a near-optimal solver rather than learning from scratch -- is the primary determinant of policy quality, while carbon-quota compliance and cost minimization occupy distinct, well-characterized points on a shared operating frontier. \S\ref{sec:limitations} states the scope and limitations of this evaluation.

\begin{thebibliography}{99}
\bibitem{sourcepaper} G.~Wang, D.~Mi, P.~He, J.~Wang, F.~Wang, M.~Fotuhi-Firuzabad, and S.~Sun, ``Deep Reinforcement Learning-Based Migration and Scheduling Strategy for Data Centers in Electricity and Carbon Markets,'' \textit{IEEE Trans. Ind. Appl.}, 2026, doi:10.1109/TIA.2026.3691718.
\bibitem{stooke2020} A.~Stooke, J.~Achiam, and P.~Abbeel, ``Responsive Safety in Reinforcement Learning by PID Lagrangian Methods,'' in \textit{Proc. 37th Int. Conf. Machine Learning (ICML)}, 2020, pp.~9133--9143.
\bibitem{radovanovic2023} A.~Radovanovi\'{c} \textit{et al.}, ``Carbon-Aware Computing for Datacenters,'' \textit{IEEE Trans. Power Syst.}, vol.~38, no.~2, pp.~1270--1280, Mar.~2023.
\bibitem{wiesner2021} P.~Wiesner, I.~Behnke, D.~Scheinert, K.~Gontarska, and L.~Thamsen, ``Let's Wait Awhile: How Temporal Workload Shifting Can Reduce Carbon Emissions in the Cloud,'' in \textit{Proc. 22nd Int. Middleware Conf. (Middleware)}, 2021, pp.~260--272.
\bibitem{bengio2021} Y.~Bengio, A.~Lodi, and A.~Prouvost, ``Machine Learning for Combinatorial Optimization: A Methodological Tour d'Horizon,'' \textit{Eur. J. Oper. Res.}, vol.~290, no.~2, pp.~405--421, 2021.
\bibitem{vinyals2015} O.~Vinyals, M.~Fortunato, and N.~Jaitly, ``Pointer Networks,'' in \textit{Advances in Neural Information Processing Systems (NeurIPS)}, 2015, pp.~2692--2700.
\bibitem{kool2019} W.~Kool, H.~van Hoof, and M.~Welling, ``Attention, Learn to Solve Routing Problems!,'' in \textit{Proc. Int. Conf. Learning Representations (ICLR)}, 2019.
\bibitem{ross2011} S.~Ross, G.~Gordon, and D.~Bagnell, ``A Reduction of Imitation Learning and Structured Prediction to No-Regret Online Learning,'' in \textit{Proc. 14th Int. Conf. Artificial Intelligence and Statistics (AISTATS)}, 2011, pp.~627--635.
\bibitem{ding2020} J.-Y.~Ding \textit{et al.}, ``Accelerating Primal Solution Findings for Mixed Integer Programs Based on Solution Prediction,'' in \textit{Proc. AAAI Conf. Artificial Intelligence}, 2020, pp.~1452--1459.
\bibitem{lee2026} H.~Lee, P.~Prabawa, D.-H.~Choi, and J.~Kim, ``Hierarchical Multi-Agent Reinforcement Learning for Carbon-Aware AI Data Centers in Power Distribution Systems,'' arXiv:2607.03324, 2026.
\bibitem{hattay2026} A.~Hattay \textit{et al.}, ``On the Role of DAG Topology in Energy-Aware Cloud Scheduling: A GNN-Based Deep Reinforcement Learning Approach,'' arXiv:2604.09202, 2026.
\bibitem{zhang2022} K.~Zhang, P.~Wang, N.~Gu, and T.~D.~Nguyen, ``GreenDRL: Managing Green Datacenters Using Deep Reinforcement Learning,'' in \textit{Proc. 13th ACM Symp. Cloud Computing (SoCC)}, 2022, pp.~445--460.
\bibitem{benmhamed2026} A.~Ben Mhamed, A.~Kamal-Idrissi, and A.~El Fallah Seghrouchni, ``Learning Branching Policies for MILPs with Proximal Policy Optimization,'' in \textit{Proc. AAAI Conf. Artificial Intelligence}, 2026.
\bibitem{relwork_branch2} ``Speeding Up Mixed-Integer Programming Solvers with Sparse Learning for Branching,'' arXiv preprint, 2026. Author attribution not independently verified at the time of writing; retained as a category reference for solver-branching-imitation work and should be confirmed or replaced before camera-ready.
\end{thebibliography}

\clearpage
\onecolumn

\begin{figure}[t]
\centering
\includegraphics[width=0.95\textwidth]{figures/fig1_architecture.png}
\caption{DFPS architecture and training pipeline: a shared bilinear scorer, batched over the entire ready-task frontier in one forward pass, combines a topology-aware neighborhood-attention encoder with a horizon-cached temporal encoder; training proceeds through solver-imitation followed by a light PID-Lagrangian PPO fine-tune.}
\label{fig:architecture}
\end{figure}

\begin{figure}[t]
\centering
\includegraphics[width=0.95\textwidth]{figures/fig2_workflow_comparison.png}
\caption{Source-paper training workflow (from-scratch on-policy PPO, full cost paid per seed) versus the proposed solver-distillation workflow (one-time solver cost amortized across all seeds).}
\label{fig:workflow}
\end{figure}

\begin{figure}[t]
\centering
\includegraphics[width=0.95\textwidth]{figures/fig3_main_benchmark.png}
\caption{Main benchmark comparison: DFPS (blue) versus 14 external benchmarks plus the minimal-distilled ablation control (orange) on the test split (100 scenarios); TAGS and HCBS-PPO (our other rejected candidates) shown in teal; TAGS-Compliant (green), the deliberately quota-compliant operating point from \S\ref{sec:carbon}, shown for reference but excluded from the formal win/tie/loss count since it targets a different (compliance-constrained) objective.}
\label{fig:main_benchmark}
\end{figure}

\begin{figure}[t]
\centering
\includegraphics[width=0.95\textwidth]{figures/fig4_generalization.png}
\caption{DFPS generalization to out-of-distribution workload counts and sensitivity to deadline-slack scale.}
\label{fig:generalization}
\end{figure}

\begin{figure}[t]
\centering
\includegraphics[width=0.95\textwidth]{figures/fig5_ablation_robustness.png}
\caption{(a) Ablation: removing the RL fine-tune phase or the topology encoder produces no measurable change; removing solver distillation (TAGS) costs 5.6\% and triples seed variance. (b) DFPS robustness to price/CEF sensor noise and bandwidth-signal dropout.}
\label{fig:ablation_robustness}
\end{figure}

\begin{figure}[t]
\centering
\includegraphics[width=0.8\textwidth]{figures/fig6_pareto.png}
\caption{Performance-efficiency Pareto: test-split system cost versus total isolated training compute across 5 seeds (log scale). DFPS (star) and Minimal-distilled (orange diamond) occupy essentially the same point -- identical training cost (both reuse the same one-time CP-SAT teacher), near-identical performance, despite Minimal-distilled having no topology or batching machinery -- and both sit closer to the strongest, most expensive cluster than either from-scratch frontier-batched candidate.}
\label{fig:pareto}
\end{figure}

\begin{figure}[t]
\centering
\includegraphics[width=0.95\textwidth]{figures/fig8_calibration.png}
\caption{Real (GEFCom2014) vs.\ synthetic day-to-day (lag-96) autocorrelation for price and carbon-intensity drivers. Our synthetic generator is considerably more temporally predictable than real data, especially for the renewable/CEF dynamics -- a quantified, disclosed limitation of the benchmark, not merely a validation exercise.}
\label{fig:calibration}
\end{figure}

\begin{figure}[t]
\centering
\includegraphics[width=0.95\textwidth]{figures/fig7_real_data_validation.png}
\caption{Real-trace-calibrated held-out generalization test (50 scenarios built from real GEFCom2014 price and wind-zone data; no retraining). DFPS's relative ranking closely replicates the synthetic-test-split result (Fig.~\ref{fig:main_benchmark}), evidence of external validity, though the additional real-world headroom above the greedy heuristic is not captured disproportionately by DFPS.}
\label{fig:real_data}
\end{figure}

\clearpage

\input{../ANALYSIS/tables/main_comparison.tex}
\input{../ANALYSIS/tables/wtl.tex}
\input{../ANALYSIS/tables/ablation.tex}
\input{../ANALYSIS/tables/efficiency.tex}

\end{document}
